\section{Implications of the geometry-first workflow}
\label{sec:geometry-first-implications}

The significance of the present construction is not that one more scalar
cosmology has been added to the literature.  The calculation demonstrates a
different route from microscopic assumptions to observations.

At the BHSM level, the starting object is a constrained geometric action.
Physical perturbations are selected only after the constraint surface,
boundary conditions, gauge quotient, and physical Hessian are specified.  At
the cosmological level, the same methodology selects a physical compact
topographic carrier, induces its metric response, transports light through
the resulting geometry, and ends with observer-frame kernels.

The resulting chain is
\begin{equation}
\boxed{
\text{geometry}
\rightarrow
\text{enclosure / physical reduction}
\rightarrow
\text{response}
\rightarrow
\text{optical propagation}
\rightarrow
\text{observation}.
}
\end{equation}

This viewpoint also clarifies the relation between BHSM and cosmology.  The
cosmological sector is not an independent effective theory appended to BHSM.
It is a large-scale test of whether the same reduction logic can generate
observable response after the physical mode count is fixed.

\subsection{Predictive content}

Several predictions or null conditions arise from the reduction itself:

\begin{enumerate}
\item the exceptional lowest scalar harmonic is not promoted to the ordinary
      propagating scalar carrier;
\item the first physical scalar carrier is the $n=2$ branch used in the
      observable transfer;
\item one scalar supernova calibration has rank one on a two-dimensional
      physical state and cannot determine the remaining state direction;
\item the pure one-mode response has a common leading observer-sky axis across
      linear observables;
\item a model with $m$ independent physical stochastic texture channels has
      topographic covariance rank at most $m$ after standard covariance and
      measurement noise are accounted for;
\item failed homogeneous and foreground reductions are rejected rather than
      repaired by post-hoc parameter changes.
\end{enumerate}

These statements differ in maturity and claim status, but they share a common
feature: they constrain future observations rather than merely describing
existing data.

\subsection{What remains open}

The present paper does not claim a complete microscopic BHSM derivation of the
cosmological source normalization.  Important upstream quantities remain
action-level open, including the normalized compact full-preimage background,
the common-domain local metric/topographic matrix elements, the exact
seam-to-topography export normalization, and the statistical law of rare
source events.

The geometry-first methodology makes these absences explicit.  They are not
replaced by fitted optical coefficients.  The photon/observable propagation
can therefore be used while the remaining upstream source calculation
continues.

\subsection{Scientific criterion}

The proposed workflow succeeds only if the same reduced physical structure
predicts multiple observables more economically and consistently than
representation-specific alternatives.  It should be rejected if future data
require incompatible axes, incompatible redshift transfer, excessive
cross-observable covariance rank, or additional independent response channels
not supplied by the retained action.

The purpose of geometry-first modeling is therefore not to avoid
falsification.  It is to move the point of falsification closer to the
underlying physical assumptions.
